\documentclass[11pt,a4paper]{article}
\usepackage[T1]{fontenc}
\usepackage[utf8]{inputenc}
\usepackage[french]{babel}
\usepackage[a4paper,margin=2.5cm]{geometry}
\usepackage{amsmath,amssymb}
\usepackage{array}
\usepackage{booktabs}
\usepackage{longtable}
\usepackage{hyperref}
\usepackage{xcolor}
\usepackage{listings}
\usepackage{tikz}
\usetikzlibrary{arrows.meta,positioning,shapes.geometric}

\hypersetup{
  colorlinks=true,
  linkcolor=blue,
  urlcolor=blue,
  pdftitle={Architecture du compilateur ifcc},
  pdfauthor={OpenAI}
}

\definecolor{codegray}{RGB}{245,245,245}
\definecolor{coderule}{RGB}{220,220,220}

\lstset{
  basicstyle=\ttfamily\small,
  backgroundcolor=\color{codegray},
  frame=single,
  rulecolor=\color{coderule},
  columns=fullflexible,
  keepspaces=true,
  showstringspaces=false,
  breaklines=true
}

\title{Architecture du compilateur \texttt{ifcc}}
\author{}
\date{}

\begin{document}
\maketitle

\section{Vue d'ensemble du pipeline}

Le compilateur \texttt{ifcc} transforme un fichier source C en fichier assembleur en plusieurs étapes successives~:

\begin{center}
\begin{tikzpicture}[
  node distance=0.9cm,
  box/.style={draw, rounded corners, align=center, minimum width=8cm, minimum height=1cm},
  arr/.style={-{Latex[length=2.5mm]}, thick}
]
  \node (c) {Fichier \texttt{.c}};
  \node[box, below=of c] (antlr) {\textbf{ANTLR4}\\Lexer + Parser $\rightarrow$ Arbre syntaxique (AST)};
  \node[box, below=of antlr] (sym) {\textbf{SymbolTableVisitor}\\Passe 1 : vérifications sémantiques};
  \node[box, below=of sym] (codegen) {\textbf{CodeGenVisitor}\\Passe 2 : construction du CFG (IR)};
  \node[box, below=of codegen] (asm) {\textbf{CFG $\rightarrow$ \texttt{gen\_asm()}}\\Émission de l'assembleur final};
  \node[below=of asm] (s) {Fichier \texttt{.s} (ARM64 ou x86-64)};

  \draw[arr] (c) -- (antlr);
  \draw[arr] (antlr) -- (sym);
  \draw[arr] (sym) -- (codegen);
  \draw[arr] (codegen) -- (asm);
  \draw[arr] (asm) -- (s);
\end{tikzpicture}
\end{center}

\section{Étape 1 --- La grammaire ANTLR4 (\texttt{ifcc.g4})}

ANTLR4 lit la grammaire et génère automatiquement un lexer et un parser en C++.

\subsection{Rôle du lexer}

Le lexer découpe le texte source en \emph{tokens}. Par exemple~:

\begin{lstlisting}
int main() { return 42; }
\end{lstlisting}

est transformé en~:

\begin{lstlisting}
[INT][ID:"main"][LPAREN][RPAREN][LBRACE][RETURN][CONST:"42"][SEMI][RBRACE]
\end{lstlisting}

\subsection{Rôle du parser}

Le parser assemble les tokens selon les règles de la grammaire et construit un AST (arbre syntaxique abstrait).

Par exemple, l'expression \texttt{3 + 4 * 2} donne une structure dans laquelle la multiplication est plus profonde que l'addition~:

\begin{center}
\begin{tikzpicture}[
  level distance=1.1cm,
  sibling distance=2.4cm,
  every node/.style={align=center}
]
\node {\texttt{add\_expr}}
  child { node {\texttt{term}} child { node {\texttt{CONST:3}} } }
  child { node {\texttt{term}}
    child { node {\texttt{CONST:4}} }
    child { node {\texttt{CONST:2}} }
  };
\end{tikzpicture}
\end{center}

La hiérarchie des règles encode les priorités opératoires~:

\begin{center}
\texttt{expr $>$ eq\_expr $>$ rel\_expr $>$ bitor\_expr $>$ xor\_expr $>$ and\_expr $>$ add\_expr $>$ term $>$ factor}
\end{center}

Plus une règle est profonde dans l'arbre, plus elle a de priorité. Ainsi, \texttt{term} (\texttt{*}, \texttt{/}, \texttt{\%}) est plus profond que \texttt{add\_expr} (\texttt{+}, \texttt{-})~: la multiplication est donc évaluée avant l'addition.

ANTLR génère dans \texttt{compiler/generated/} les fichiers suivants~:
\begin{lstlisting}
ifccParser.cpp
ifccLexer.cpp
ifccBaseVisitor.cpp
\end{lstlisting}

Ces fichiers sont générés automatiquement et ne doivent jamais être modifiés à la main.

\section{Étape 2 --- Le \texttt{SymbolTableVisitor} (passe sémantique)}

Cette première visite de l'AST ne produit pas de code~: elle vérifie uniquement que le programme est sémantiquement correct.

\subsection{Le pattern Visitor}

ANTLR génère une classe \texttt{ifccBaseVisitor} contenant une méthode \texttt{visitXxx} pour chaque règle de grammaire. On hérite de cette classe et on redéfinit uniquement les méthodes utiles. Lorsqu'on appelle \texttt{visit(noeud)}, ANTLR choisit automatiquement la bonne méthode.

\subsection{Pile de scopes}

Le \texttt{SymbolTableVisitor} maintient une pile de scopes~:

\begin{lstlisting}
scopes_ = [
  {"x": true, "y": false},   // scope fonction
  {"i": true}                // scope bloc interne
]
\end{lstlisting}

Chaque scope est une \texttt{map<string,bool>} associant à chaque variable un booléen indiquant si elle a été utilisée.

Les opérations principales sont~:

\begin{itemize}
  \item \texttt{enterScope()} : empile une nouvelle map vide ;
  \item \texttt{exitScope()} : dépile le scope courant et émet un warning pour chaque variable jamais utilisée ;
  \item \texttt{declareVar(name)} : vérifie qu'il n'existe pas déjà de variable du même nom dans le scope courant, puis l'ajoute ;
  \item \texttt{isVisible(name)} : cherche une variable du scope le plus interne vers le plus externe ;
  \item \texttt{markUsed(name)} : marque comme utilisée la déclaration visible la plus interne.
\end{itemize}

\subsection{Table des fonctions}

Le visiteur maintient aussi une table des fonctions~:

\begin{lstlisting}
map<string, FuncInfo>
\end{lstlisting}

Elle est remplie en deux sous-passes~:

\begin{enumerate}
  \item \texttt{collectSignatures()} : enregistre pour chaque fonction son nom, son nombre de paramètres et son type de retour ;
  \item visite des corps de fonctions : chaque appel \texttt{f(a, b)} est vérifié par rapport à la table collectée.
\end{enumerate}

Cette organisation permet les appels avant définition (\emph{forward references}).

\subsection{Erreurs et warnings}

Les cas suivants provoquent \texttt{hasError = true} (sortie en erreur, pas d'assembleur)~:

\begin{itemize}
  \item variable utilisée avant déclaration ;
  \item variable déclarée deux fois dans le même scope ;
  \item appel à une fonction non définie ;
  \item mauvais nombre d'arguments ;
  \item fonction définie deux fois.
\end{itemize}

Les cas suivants déclenchent seulement un warning~:

\begin{itemize}
  \item variable déclarée mais jamais utilisée.
\end{itemize}

\section{Étape 3 --- La représentation intermédiaire (IR)}

Avant d'émettre l'assembleur, \texttt{ifcc} passe par une représentation intermédiaire indépendante de l'architecture cible.

\subsection{Classe \texttt{IRInstr}}

Une instruction IR représente une opération atomique~:

\begin{lstlisting}
enum Operation {
  ldconst, copy, add, sub, mul, div, mod,
  cmp_eq, cmp_ne, cmp_lt, cmp_gt,
  bit_or, bit_xor, bit_and, not_op, call, ...
};
\end{lstlisting}

Chaque instruction contient~:
\begin{itemize}
  \item une opération ;
  \item un type ;
  \item une liste de paramètres correspondant à des variables IR.
\end{itemize}

Exemples~:

\begin{lstlisting}
add  INT  ["!tmp0", "x", "y"]    // !tmp0 = x + y
copy INT  ["!retval", "!tmp0"]    // !retval = !tmp0
\end{lstlisting}

Toutes les variables IR appartiennent à un espace de noms plat. Chaque variable possède un slot sur la pile, identifié par un \texttt{SymbolIndex}. Les temporaires prennent des noms comme \texttt{!tmp0}, \texttt{!tmp1}, etc., et la valeur de retour s'appelle \texttt{!retval}.

\subsection{Classe \texttt{BasicBlock}}

Un bloc de base est une séquence d'instructions sans branchement interne. Il se termine par une décision de contrôle de flot~:

\begin{lstlisting}
BasicBlock* exit_true;
BasicBlock* exit_false;
std::string test_var_name;
bool isReturnExit;
\end{lstlisting}

Interprétation~:
\begin{itemize}
  \item \texttt{exit\_true == nullptr} et \texttt{isReturnExit} : le bloc produit l'épilogue puis \texttt{ret} ;
  \item \texttt{exit\_false == nullptr} : saut inconditionnel vers \texttt{exit\_true} ;
  \item les deux pointeurs non nuls : branchement conditionnel sur \texttt{test\_var\_name}.
\end{itemize}

\subsection{Classe \texttt{CFG}}

Chaque fonction est représentée par un graphe de flot de contrôle (\emph{Control Flow Graph})~:

\begin{lstlisting}
CFG "main"
├── BasicBlock .Lmain_0   (entrée)
├── BasicBlock .Lmain_1   (branche then)
├── BasicBlock .Lmain_2   (merge)
└── ...
\end{lstlisting}

Le \texttt{CFG} gère également~:
\begin{itemize}
  \item la table des symboles physiques, par exemple \texttt{SymbolIndex["x"] = 4} ;
  \item la taille de la frame de pile, alignée sur 16 octets ;
  \item le prologue et l'épilogue de fonction ;
  \item le passage des paramètres.
\end{itemize}

\section{Étape 4 --- Le \texttt{CodeGenVisitor} (construction du CFG)}

La seconde visite de l'AST ne vérifie plus le programme~: elle construit les \texttt{CFG} des fonctions.

\subsection{Convention de retour des visiteurs d'expressions}

Chaque méthode \texttt{visitXxx} correspondant à une expression retourne un \texttt{antlrcpp::Any} contenant le nom de la variable IR où se trouve le résultat.

Une fonction utilitaire permet de récupérer cette valeur~:

\begin{lstlisting}
std::string evalExpr(ParseTree* node) {
    return std::any_cast<std::string>(visit(node));
}
\end{lstlisting}

Exemple pour \texttt{3 + x}~:

\begin{lstlisting}
visitAdd_expr
  -> evalExpr(term "3")
     -> visitFactor -> ldconst !tmp0, 3 -> retourne "!tmp0"
  -> evalExpr(term "x")
     -> visitFactor -> retourne "x"
  -> add !tmp1, !tmp0, x
  -> retourne "!tmp1"
\end{lstlisting}

\subsection{Gestion du shadowing}

Le langage C autorise qu'une variable interne masque une variable externe du même nom. Or le \texttt{CFG} utilise un espace de noms plat. La solution consiste à maintenir une pile de renommage~:

\begin{lstlisting}
nameScopes_ : pile de maps nom_logique -> nom_IR
\end{lstlisting}

Exemple~:
\begin{itemize}
  \item \texttt{int x} dans le scope externe $\rightarrow$ nom IR \texttt{"x"} ;
  \item \texttt{int x} dans un scope interne $\rightarrow$ nom IR \texttt{"x@0"}.
\end{itemize}

La fonction \texttt{resolveVar("x")} cherche du scope le plus interne vers le plus externe et retourne le bon nom IR.

\subsection{Gestion du flot de contrôle}

Pour une structure~:

\begin{lstlisting}
if (cond) { then } else { else }
\end{lstlisting}

le \texttt{CodeGenVisitor} construit~:

\begin{lstlisting}
current_bb:
  [évalue cond -> condVar]
  test_var_name = condVar
  exit_true  = thenBB
  exit_false = elseBB

thenBB:
  [visite then body]
  exit_true = mergeBB   // si pas de return

elseBB:
  [visite else body]
  exit_true = mergeBB   // si pas de return

mergeBB:
  [suite du code]
\end{lstlisting}

Le drapeau \texttt{isReturnExit} empêche de raccorder un bloc terminé par \texttt{return} au bloc de fusion. Lorsqu'un \texttt{return} est rencontré, le bloc courant est marqué comme bloc de sortie et un \emph{sink block} absorbe le code mort éventuel qui suivrait.

\section{Étape 5 --- Émission de l'assembleur (\texttt{gen\_asm})}

La méthode \texttt{CFG::gen\_asm()} effectue les opérations suivantes~:

\begin{enumerate}
  \item émission du prologue de fonction ;
  \item génération du code de chaque \texttt{BasicBlock} ;
  \item émission des branchements et de l'épilogue.
\end{enumerate}

Chaque \texttt{IRInstr::gen\_asm()} traduit une instruction IR en assembleur concret. Deux variantes existent selon l'architecture cible.

\subsection{Exemple ARM64}

\begin{lstlisting}
// ARM64 (Apple Silicon)
case add:
    ldr w8, [x29, #-offset_gauche]
    ldr w9, [x29, #-offset_droite]
    add w8, w8, w9
    str w8, [x29, #-offset_dest]
\end{lstlisting}

\subsection{Exemple x86-64}

\begin{lstlisting}
// x86-64 (Linux/Intel)
case add:
    movl -offset_gauche(%rbp), %eax
    addl -offset_droite(%rbp), %eax
    movl %eax, -offset_dest(%rbp)
\end{lstlisting}

Dans cette architecture pédagogique, toutes les variables résident en mémoire, sur la pile. Il n'y a pas d'allocation de registres~: cela simplifie fortement le compilateur, au prix d'un code moins optimisé.

\section{Exemple complet}

Considérons le programme suivant~:

\begin{lstlisting}
int main() {
    int x;
    x = 3 + 4;
    return x;
}
\end{lstlisting}

\subsection{Parsing}

Le parser construit un AST comportant une définition de fonction, puis une séquence d'instructions~: déclaration, affectation, retour.

\subsection{Passe sémantique}

Le \texttt{SymbolTableVisitor} effectue les actions suivantes~:

\begin{itemize}
  \item création du scope de \texttt{main} ;
  \item déclaration de \texttt{x} avec \texttt{used = false} ;
  \item affectation à \texttt{x} $\Rightarrow$ \texttt{x} est marquée utilisée ;
  \item lecture de \texttt{x} dans \texttt{return} $\Rightarrow$ \texttt{x} reste marquée utilisée ;
  \item sortie du scope sans warning.
\end{itemize}

\subsection{Génération IR}

Le \texttt{CodeGenVisitor} construit le \texttt{CFG} de \texttt{main}~:

\begin{lstlisting}
CFG "main" créé
!retval alloué
.Lmain_0 créé comme bloc courant
int x -> add_to_symbol_table("x") -> offset 4

x = 3 + 4 :
  ldconst !tmp0, 3
  ldconst !tmp1, 4
  add     !tmp2, !tmp0, !tmp1
  copy    x, !tmp2

return x :
  copy    !retval, x
  isReturnExit = true
  sink BB créé
\end{lstlisting}

\subsection{Émission finale}

L'assembleur final contient typiquement~:

\begin{itemize}
  \item un prologue (sauvegarde de \texttt{x29/x30} sur ARM64 ou \texttt{rbp} sur x86-64) ;
  \item les instructions correspondant aux opérations IR ;
  \item un épilogue chargeant \texttt{!retval} dans le registre de retour ;
  \item l'instruction \texttt{ret}.
\end{itemize}

\section{Fonctions et absence d'imbrication}

Le compilateur \texttt{ifcc} ne supporte pas les fonctions imbriquées (une fonction définie à l'intérieur d'une autre). Cette contrainte est imposée directement par la grammaire ANTLR4~:

\begin{lstlisting}
prog : func_def+ ;
func_def : rettype ID '(' params? ')' '{' stmt* '}' ;
stmt : decl_stmt | expr_stmt | return_stmt | block_stmt | if_stmt | while_stmt ;
// pas de func_def dans stmt !
\end{lstlisting}

Ainsi, le programme suivant est invalide~:

\begin{lstlisting}
int main() {
    int helper() { return 1; }  // interdit
    return helper();
}
\end{lstlisting}

Les fonctions doivent être définies au même niveau~:

\begin{lstlisting}
int helper() { return 1; }
int main()   { return helper(); }
\end{lstlisting}

\subsection{Comportement du \texttt{SymbolTableVisitor}}

Dans \texttt{visitProg}, deux passes sont effectuées~:

\begin{lstlisting}
collectSignatures(ctx);
for (auto* f : ctx->func_def())
    this->visit(f);
\end{lstlisting}

Dans \texttt{visitFunc\_def}, la pile de scopes est complètement réinitialisée pour chaque fonction~:

\begin{lstlisting}
scopes_.clear();
enterScope();
\end{lstlisting}

Chaque fonction possède donc un environnement totalement indépendant.

\subsection{Appels de fonctions}

Lorsqu'un appel comme \texttt{helper()} est rencontré dans \texttt{main}, le compilateur ne visite pas le corps de \texttt{helper}. Il vérifie uniquement la signature~:

\begin{lstlisting}
auto it = functions_.find("helper");
// vérifie existence + nombre d'arguments
\end{lstlisting}

La table \texttt{functions\_} est globale, contrairement aux scopes qui sont locaux à chaque fonction.

\subsection{Résumé}

\begin{lstlisting}
visitProg
├── collectSignatures → functions_ = {"helper", "main"}
│
├── visitFunc_def "helper"
│   ├── scopes_.clear()
│   ├── enterScope()
│   └── exitScope()
│
└── visitFunc_def "main"
    ├── scopes_.clear()
    ├── enterScope()
    ├── appel helper() → vérification dans functions_
    └── exitScope()
\end{lstlisting}

Les fonctions imbriquées nécessiteraient un mécanisme de capture d'environnement (comme les closures), ce qui dépasse le modèle de \texttt{ifcc} et même du langage C standard.

\section{Fonctionnalités du compilateur}

\subsection{Type entier 32 bits}

Le type entier est défini dans \texttt{Type.h}. La fonction \texttt{typeSize(INT)} retourne 4.

\begin{itemize}
  \item Sur ARM64 : registres \texttt{w} (32 bits)
  \item Sur x86-64 : suffixe \texttt{l} (\texttt{movl}, \texttt{addl})
\end{itemize}

Toute l'IR travaille en 32 bits.

\subsection{Constantes}

\begin{itemize}
  \item Entières : supportées
  \item Caractère : non supportées
\end{itemize}

\begin{lstlisting}
CONST : [0-9]+ ;
\end{lstlisting}

\subsection{Opérateurs}

\begin{itemize}
  \item +, -, *, /, \% → IRInstr arithmétiques
  \item |, \&, ^ → bitwise
  \item ==, !=, <, > → comparaisons
\end{itemize}

\subsection{Affectation}

Retourne une valeur, permet le chaînage.

\subsection{Fonctions}

\begin{itemize}
  \item Paramètres copiés sur la pile
  \item Types : int / void
\end{itemize}

\subsection{Contrôle de flux}

\begin{itemize}
  \item if / else
  \item while
\end{itemize}

Implémentés via des BasicBlocks.

\subsection{Scopes et shadowing}

Renommage des variables internes (x@0, x@1).

\subsection{Vérifications sémantiques}

\begin{itemize}
  \item déclaration avant usage
  \item pas de doublon
  \item warning si non utilisée
\end{itemize}

\section{Conclusion}

L'architecture de \texttt{ifcc} repose sur une séparation claire entre~:

\begin{itemize}
  \item l'analyse syntaxique (ANTLR4) ;
  \item la vérification sémantique (\texttt{SymbolTableVisitor}) ;
  \item la construction d'une représentation intermédiaire structurée (\texttt{CodeGenVisitor}, \texttt{IRInstr}, \texttt{BasicBlock}, \texttt{CFG}) ;
  \item l'émission finale de l'assembleur.
\end{itemize}

Cette organisation rend le compilateur plus modulaire, plus lisible et plus facile à faire évoluer, par exemple pour ajouter de nouvelles constructions du langage ou améliorer les optimisations.

\end{document}
